\documentclass[12pt,a4paper,onecolumn]{article}
\title{Digest of The New York Times}
\author{The New York Times}
% \definecolor{dark-green}{rgb}{0,0.25,0}
% \usepackage[utf8]{inputenc}

\usepackage{fontspec,xunicode,xltxtra}
\setmainfont[Ligatures=Required,Mapping=tex-text]{Georgia}
% \setmainfont[Ligatures=Required,Mapping=tex-text]{Palatino Linotype}
\setsansfont[Ligatures=Required,Mapping=tex-text]{Myriad Pro}
\setmonofont[Ligatures=Required,Mapping=tex-text]{Courier New}
\usepackage{xeCJK}
\setCJKmainfont{宋体}
\setCJKmonofont{黑体}

% \usepackage{indentfirst}


\usepackage[dvips,dvipsnames,svgnames]{xcolor}
\definecolor{light-gray}{gray}{0.95}

%生成PDF的链接
\usepackage[dvipdfmx]{hyperref}
\hypersetup{
    bookmarks=true,         % show bookmarks bar?
    bookmarksopen=true,
    pdfpagemode=UseNone,    % options: UseNode, UseThumbs, UseOutlines, FullScreen
    pdfstartview=FitH,
    pdfborder=1,
    pdfhighlight=/P,
    pdfauthor={wuxch},
    unicode=true,           % non-Latin characters in Acrobat’s bookmarks
    colorlinks,             % false: boxed links; true: colored links
    linkcolor=blue,         % color of internal links
    citecolor=blue,        % color of links to bibliography
    filecolor=magenta,      % color of file links
    urlcolor=cyan           % color of external links
}
% 使用目录
\usepackage{makeidx}
% 使用图形
\usepackage{graphicx}


\usepackage{titletoc}
\dottedcontents{section}[1.5em]{}{2.3em}{1pc}

% % 页眉页脚的设定
% \usepackage{fancyhdr}
% \pagestyle{fancy}
% \lhead{\includegraphics[height=0.2in]{zte_logo}}

% \usepackage{footmisc}
\usepackage{pfnote}
\usepackage{fnpos}

% % 嵌入的代码显示
% \usepackage{listings}
% \lstset{language=C++, breaklines, extendedchars=false}
% \lstset{keywordstyle=\color{blue},
%         commentstyle=\color{SeaGreen},
%         stringstyle=\ttfamily,
%         showstringspaces=false,
%         backgroundcolor=\color{light-gray}}
% \lstset{basicstyle=\small}

\usepackage{booktabs}

\usepackage{tipa}
\usepackage{phonetic}
\usepackage[gen]{eurosym}

% \usepackage{CJK}
\usepackage{makeidx}
% \usepackage{glossary}

\usepackage{fancyhdr}

\usepackage{lastpage}
\cfoot{\textsf{\thepage\ of \pageref*{LastPage}}}

\pagestyle{fancy}
\fancyhead{}
\renewcommand{\headrulewidth}{0pt}
\pagenumbering{arabic}

\usepackage{graphicx}
\usepackage{picinpar}

\setlength{\parindent}{0in}
\setlength{\parskip}{2.5ex plus4mm minus3mm}
% \setlength{\parskip}{.3\baselineskip}

\setcounter{section}{0}

% bibtex
\usepackage[super]{natbib}
% \usepackage[globalcitecopy,labelstoglobalaux]{bibunits}
% \makeatletter
% \renewcommand{\@bibunitname}{\jobname.\the\@bibunitauxcnt}
% \makeatother

% \renewcommand{\refname}{缩略语}
% \renewcommand{\refname}{}

\makeatletter                   % Reference list option change
\renewcommand\@biblabel[1]{#1:  } % from [1] to 1: 
\makeatother                    %

\newcommand{\citenumfont}[1]{\textsf{#1}}

% \usepackage{tipa}
% \usepackage{phonetic}

\newcommand{\phonetic}[1]{{\fontspec[Colour=CC00CC]{Microsoft Sans Serif} #1}}


%%% Local Variables:
%%% TeX-master: "medical_trivia"
%%% End:

\makeindex
\begin{document}
\date{}
\thispagestyle{empty}
\renewcommand\contentsname{\textsf{Digest of The New York Times}}
\textsf{\tableofcontents}
\pagebreak
\setcounter{page}{1}
\section{Mini Versions of Big-Box Stores}

BOTHELL, Wash. -- During the current economic downturn, as many companies are closing stores and
cutting costs, it might seem counterintuitive to be opening new stores.

Not here in Bothell, around 20 miles northeast of Seattle, where in January, OfficeMax opened one of
its three new concept stores in the Seattle area that offer a pared-down selection of its most
popular products. Each of the new stores, called Ink Paper Scissors, covers only 2,000 square feet
-- about a ninth the size of a typical OfficeMax -- and offers basics like copy-making supplies and
printer-cartridge refills.

Retailers like OfficeMax are opening scaled-down versions of their stores or inventing outlets
entirely to test new concepts without a hefty investment. The stores are a relatively safe bet
despite the recession because the space is cheaper and the stores require less inventory, fewer
employees and smaller spaces.

OfficeMax is not the only retailer giving new concept stores a try. Most are significantly smaller
than their typical stores and focus on one set of products. Last year, for example, Wal-Mart opened
four specialty food stores in the Phoenix area, RadioShack unveiled three high-end wireless shops in
Dallas, and Best Buy created 30 mobile phone stores.

``If you've got the wherewithal, everyone is thinking about smaller sizes,'' said Lee Peterson, vice
president for brand and creative services at WD Partners in Columbus, Ohio, which has helped
retailers design these stores.

More small-format and new-concept stores are likely to be on the way as retailers try to lure
customers back, according to a survey of retailers, manufacturers and consultants. Nearly 46 percent
of the respondents said they expected the number of formats to increase in the next five years,
according to a survey in February by Dechert-Hampe \& Company, a marketplace management consultancy.

The stores are opening even as some companies are declaring bankruptcy, closing stores and reporting
double-digit earnings drops. As the economy contracts and consumers tighten their spending, more
closings are likely. In March, the most recent data available, Americans were saving 4.2 percent of
their income, after taxes, up from 0.2 percent a year earlier, according to the Commerce Department.

``Obviously, it's not an optimal time,'' said Ryan Vero, OfficeMax's chief merchandising officer.
``But this makes for a great test -- it can't be any worse.''

Smaller formats also allow companies to enter new markets in urban or rural areas that they had
bypassed during the boom. They can get into a market, test a new concept and get out quickly if it
doesn't work.

The Lowe's Companies, the home repair giant, has been hit hard by the housing downturn and has cut
its new store openings to 60 to 70 this year, from 115 last year. But the company, which is based in
Mooresville, N.C., is also experimenting with scaled-back stores in new markets. Lowe's has slowed
its plans to open stores in the hard-hit states of Florida, Arizona and California, and is aiming at
the Midwest and rural communities far from its warehouse stores.

Last year, Lowe's opened two scaled-down versions of 66,000 square feet and 80,000 square feet. An
average Lowe's store is 117,000 square feet.

Many retailers don't want to cease opening new stores altogether during the downturn, because doing
so could harm future earnings. It can take up to three years to develop and open a store. But
opening during a recession can position a retailer for success when the economy turns.

``Retailers have pulled back the reins somewhat, but they're not going to pull back entirely,'' said
Daniel Butler, vice president of retail operations at the National Retail Federation in Washington.

Smaller stores are also cheaper and less risky. An average Lowe's costs \$20 million to \$22 million
to build, not including inventory. Smaller stores save the chain an average of almost 10 percent, or
\$1.9 million, per store.

They also can attract new customers who might be put off by larger stores or consumers who shop
mostly online. Downsized or concept stores are more convenient and take less time to visit than a
large store. Lines are typically shorter, and the shopping aisles can be easier to navigate.

In some ways, retailers are going back to their roots, evoking the corner store. At many new stores,
personalized service is being emphasized, like explaining the features of a product.

``Consumers want stores that are more convenient, less time-consuming and more personal,'' said Ben
Ball, senior vice president at Dechert-Hampe \& Company. ``There is such a thing as too much
variety.''

Best Buy is aiming to lure people away from their computers and into their new Best Buy Mobile
stores with a selection of 90 cellphones and service plans from nine carriers.

Most of the stand-alone shops are 900 to 1,200 square feet, versus 40,000 square feet for a big-box
Best Buy. Some 3,000-square-foot versions have also opened. Best Buy now has 40 stand-alone stores
and plans more.

Best Buy is focusing on the cellphone and smartphone market because the segment is expected to keep
growing and the company needs to shore up flagging earnings. Best Buy's goal is to have 10 percent
of the country's cellphone market, up from about 3 percent now.

``We've been in the mobile business for years, but despite our best efforts, we haven't been able to
make much impact in cellphones,'' said Scott Moore, vice president of marketing for Best Buy Mobile.

RadioShack is also eyeing the wireless market. In December, the company opened three Point Mobl
stores in the Dallas area; they carry smartphones, among other electronics items. Each store is
1,500 square feet, or about the size of a typical RadioShack. The new stores do not mention the
parent company.

Retailing professionals doubt the experimentation with store formats is over. ``It's a very liquid
time for retailers,'' said Mr.~Peterson at WD Partners. `` 'Let's try it' is now the mantra.''

\section{Search for Supreme Court Justice Reaches Interview Stage}

WASHINGTON -- President Obama has interviewed his first prospective Supreme Court candidate, sitting
down privately in the White House for a conversation with Judge Diane P.~Wood, an official confirmed
Wednesday.

Judge Wood, who was appointed by President Bill Clinton 14 years ago to the United States Court of
Appeals for the Seventh Circuit, was attending a law conference at Georgetown University on
Wednesday. She declined to comment to reporters about any other stops on her trip to Washington from
Chicago.

The announcement of a selection for the court is not imminent, the official said, and other
finalists are expected to be interviewed. White House aides are preparing for an announcement no
sooner than late next week.

Judge Wood, 58, is well known to the president and some of his advisers because of her Chicago
connection. When she was appointed to the appeals court, she was a professor at University of
Chicago Law School, where Mr.~Obama taught constitutional law.

The meeting with Judge Wood was believed to be the president's first one-on-one interview with a
possible candidate. The official who confirmed the meeting spoke on the condition of anonymity
because of the secrecy of the selection process to fill the seat of Justice David H.~Souter, who is
retiring.

The selection process is awash in secrecy, with only a tight circle of advisers familiar with
Mr.~Obama's thinking on the matter. The president's travel schedule could provide some insight into
the timing of the announcement.

Mr.~Obama is set to spend the weekend at Camp David, attend a Memorial Day event at Arlington
National Cemetery and on Tuesday begin a three-day swing to the West Coast for appearances in Nevada
and California.

The announcement is expected to take place at the White House, where Mr.~Obama will be on Friday and
the first two days of the next week before leaving for a trip to Egypt and Europe.

The conference of legal scholars and judges that brought Judge Wood to town honored former Justice
Sandra Day O'Connor.

Judge Wood was not the only potential Supreme Court candidate in the room.

A speaker at the conference was Elena Kagan, the United States solicitor general, who is also
believed to be a candidate for the court vacancy.

Gregory B.~Craig, the White House counsel and one of the president's primary advisers in the search
for a new justice, also spoke at the conference, which was devoted to a discussion of ``Fair and
Independent Courts in a New Era.''

\section{Keeping Score on Obama vs. Netanyahu}

WASHINGTON -- After the much anticipated White House meeting on Monday between President Obama and
Prime Minister Benjamin Netanyahu of Israel, one question being asked in diplomatic circles is this:
Did Mr.~Obama give up more than he got?

The meeting between the two, their first as leaders, was mainly an exercise in breaking the ice. But
at the early stages of a relationship between the nations' leaders that is likely to be more
strained than it was during the Bush years, their dealings are being analyzed for signs of who has
the upper hand.

American and Israeli officials had predicted an exchange with some give-and-take: Mr.~Netanyahu
would try to extract from Mr.~Obama a timetable for dealing with Iran, with a deadline for Tehran to
stop enriching uranium or face serious repercussions.

In return, Mr.~Obama would push Mr.~Netanyahu to move swiftly on a peace plan with the Palestinians,
as well as to freeze the building of Jewish settlements in the West Bank.

Mr.~Netanyahu got his timetable. ``We're not going to have talks forever,'' Mr.~Obama said of Iran,
assuring Mr.~Netanyahu that he expected to know by the end of the year whether Iran was making ``a
good-faith effort to resolve differences.''

But Mr.~Obama did not get his settlement freeze. In fact, Mr.~Netanyahu told him it would be
politically difficult for him to halt the construction of settlements. That is a hurdle to the
administration's broader peace objectives because Israel's Arab neighbors have characterized a
freeze as a precondition for them to establish normal relations.

Nor did Mr.~Obama get much from Mr.~Netanyahu on a peace plan beyond his promise to make good on a
few commitments that Israel had already agreed to on the ``road map,'' an outline of peace steps
that has not gotten either Palestinians or Israelis any closer to peace since President George
W.~Bush first announced it in 2003.

Mr.~Netanyahu did agree to resume talks with Palestinians without preconditions. But he would not
explicitly endorse the notion of an eventual Palestinian state, something his predecessor, Ehud
Olmert, had already done.

``This is why I'm asking the question, did our president get suckered?'' said Martin S.~Indyk, a
former United States ambassador to Israel and director of the Saban Center at the Brookings
Institution. ``We don't know the answer yet, but unless he got something more from Bibi in that
meeting than they're telling us, that question can be asked.''

The two leaders set up working groups to deal with Iran, the Palestinian issue and Israel's Arab
neighbors. The groups will meet periodically, Israeli and American officials said. Agreeing to meet
with Israel regularly to discuss the administration's progress with Tehran keeps the pressure
squarely on the United States, analysts said.

Mr.~Netanyahu said Wednesday in Jerusalem that he was also willing to open talks with Syria, if
Syria did not make preconditions.

In offering their assessments of the meeting in Washington, a range of American and Israeli
officials generally insisted on anonymity to address the question of which leader, if either, had
gained the upper hand.

White House officials maintained Wednesday that Mr.~Obama had emerged from the meeting with his
peace goals intact. ``The president was clear, both publicly and privately, that all parties,
including the Israelis, have obligations as they relate to settlements, as they relate to Gaza, and
as they relate to two states,'' a senior White House official said.

Mr.~Obama, the official said, pressed Mr.~Netanyahu on settlements, and added that American
officials would monitor the issue closely in the next months. The official characterized Monday's
meeting as a start: ``This was the alpha meeting. It wasn't the alpha and omega meeting.''

Another administration official said Mr.~Obama's timetable on Iran was not predicated on a quid pro
quo from Israel. Mr.~Netanyahu, he noted, endorsed the diplomatic overture to Iran -- a significant
gesture, given Israel's deep fear of Iranian intentions.

Still, Mr.~Netanyahu's silence on settlements troubles lawmakers like Senator John Kerry, the
Massachusetts Democrat who is chairman of the Senate Foreign Relations Committee. Mr.~Kerry raised
the issue on Wednesday with Secretary of State Hillary Rodham Clinton when she testified before the
panel.

``It is clear that settlement activity has to cease,'' Mrs.~Clinton said, noting that she had
pressed Mr.~Netanyahu about settlements over dinner on Monday.

Israeli officials also disputed the notion that Mr.~Netanyahu had somehow outmaneuvered Mr.~Obama.
``Obama may be slightly less experienced than Netanyahu, but Obama knows exactly everything that the
U.S.~is doing,'' an Israeli official said.

The Israeli official said Mr.~Netanyahu's refusal to explicitly endorse Palestinian statehood now
did not mean he would not later.

\section{Tiananmen Now Seems Distant to China's Students}

BEIJING -- On April 30, the cellphones of the 32,630 students at Peking University, a genteel
institution widely regarded as one of China's top universities, buzzed with a text message from the
school administration. It warned students to ``pay attention to your speech and behavior'' on Youth
Day because of a ``particularly complex'' situation.

Few students had to puzzle over the meaning. Youth Day, on May 4, commemorates a 1919 student
protest against foreign imperialism and China's weakness in resisting it. Seventy years later, in
1989, students from Peking University were again massing in the center of Beijing, demanding
democracy. The student movement shook the ruling Communist Party to its core and ended with a
military crackdown and hundreds of deaths.

And if a student today proposed a pro-democracy protest?

``People would think he was insane,'' said one Peking University history major in a recent
interview. ``You know where the line is drawn. You can think, maybe talk, think about the events of
1989. You just cannot do something that will have any public influence. Everybody knows that.''

Most students also appear to accept it. For 20 years, China's government has made it abundantly
clear that students and professors should stick to the books and stay out of the streets. Students
today describe 1989 as almost a historical blip, a moment too extreme and traumatic ever to repeat.

But whether democracy still inspires them is a more complex question.

Interviews with students and teachers at Peking University, as with experts on China here and
abroad, draw a layered portrait of today's students: disinclined to protest, but also lacking the
economic grievances that helped ignite protests in 1989; proud of China's achievements and flocking
to the Communist Party, but seldom driven by ideology.

They are disturbed by government corruption and censorship and are eager to study in the West,
especially the United States. And despite the government's attempts to wipe the 1989 protests from
Chinese history, some have learned what happened. All but one of eight Peking University students
interviewed for this article, for instance, said they had managed to download an acclaimed -- and
banned -- documentary on the Tiananmen protests and view it in their dorm rooms.

``There is a stereotypical view that students are not interested in democracy. I don't buy it,''
Cheng Li, research director of the China Center at the Brookings Institution, said in an interview.
``At the very least, they have a mixed opinion of the Communist Party.''

Xia Yeliang, a Peking University professor, said many students supported democracy in theory but did
not want to risk their futures to fight for it. Students joke that they will get involved once
pro-democracy forces gather steam, he said. ``A rather high percentage of students are not
interested in politics,'' he said. ``They say, 'We know this is a good thing, but what relation does
it have to us?' They think about their personal affairs, how to get a job, how to go abroad.''

Even the mouthpiece of the Communist Party, People's Daily, laments a general lack of idealism on
campus. ``Many university students are clearly very utilitarian in their thinking,'' People's Forum,
a magazine published by People's Daily, complained this month after a conducting a student survey.
``Everything is based on 'whether or not it is useful to me,' '' the magazine said.

In fact, today's students have more to lose than did protesters 20 years ago. Then, university
students believed that their futures were endangered by a soaring inflation rate of 28 percent,
rampant government corruption and shrinking job prospects, according to a 2001 book on the Tiananmen
movement by Dingxin Zhao, a University of Chicago sociology professor. Many had lost hope in the
government's economic reforms.

Today, even students who criticize Communist rule are gratified by China's great strides.
``Sometimes we don't like the policies of our government,'' said Wang Yongli, a fourth-year physics
major. ``But on the other hand, nowadays we are proud of the country and the government because they
have moved so many people to a better life.''

The Communist Party is careful to cultivate this image, while seeking to defuse longings for
democracy by vowing to govern ``democratically.''

Officials say they oppose Western-style multiparty democracy as wrong for China, but embrace the
idea of consultation, public review and balloting under party rule. China will open up the political
system, step by step, as the country becomes wealthier and more stable, officials promise.

Some China analysts suggest that student discontent could rise if the current economic crisis clouds
their futures. China sends nine times as many students to institutions of higher education now as it
did in 1989, and competition for good jobs is fierce. Nearly one in four graduates last year could
not find work, Xinhua, the state-run news agency, reported.

But since 1989, Communist Party leaders have realized that they ignore youth at their peril. The
government is now trying to ease job anxieties with training programs and incentives for graduates
to work in rural areas. ``If you are worried, then I am more worried than you,'' Prime Minister Wen
Jiabao told one student group in December.

The party has also ratcheted up recruitment and political education, making college students the
party's fastest-growing segment, said Susan L.~Shirk, a political science professor at the
University of California, San Diego. More than 8 percent of all students were party members in 2007,
compared with fewer than 1 percent in 1989. At elite institutions like Peking University,
percentages are much higher.

Some of those students echo the party's line that Western-style democracy does not suit China.
``China has a large population, and education has a long way to go,'' said Song Chao, a Peking
University ecology major. ``Considering that, we need to put some regulations on people. The major
task for China now is development.''

Others hope to nudge the party toward reform. ``Of course, if we could become a democratic society,
we would like that,'' said another history major and party aspirant. ``But this is not something you
can achieve by radical means. What if there is chaos?''

But a majority of students seek party membership not as an ideological statement but rather as a
means to a better job, the survey published by People's Forum concluded. At Peking University, many
students say they nap through the university's much mocked, though mandatory, political thought
classes. ``Even the teachers know they are teaching rubbish,'' one senior said.

Most students will make such statements only anonymously because government control of campus speech
remains tight. Professors say some students are assigned to report to administrators if they hear
teachers adopting antigovernment lines. Most students interviewed for this article did not want to
be identified, saying their comments might be negatively noted in their files.

Five years ago, the university shut down a computer bulletin board -- a vibrant hub of information
for 300,000 users -- after the central government's education minister complained that it did not
always reflect ``the right view.'' Students say they are careful about what they write on the new,
restricted and monitored board because their identities can be traced.

Surveys show that four of five university students still rely on China's heavily censored media for
their news. But in a digital age when nearly 70,000 Chinese students are studying in the United
States and roughly 163,000 foreign students study at Chinese universities, walls against information
are porous.

One senior recalled an excruciating roundtable discussion with foreign journalists who visited
Peking University in 2007 and asked about the government crackdown on student demonstrators in 1989.
``They always ask about this June 4 incident, and we just keep silent,'' she said. ``It is not
because we don't want to talk. It is because we have no idea what exactly happened!

``I felt a little bit humiliated because we don't know our own history,'' she said. ``So I went to
the library and I read about June 4. Basically, everything was written by foreign journalists.''

The curbs on public debate can reduce even political controversies on campus to the status of
rumors. Two Peking University professors were among the first to sign Charter 08, an online
pro-democracy manifesto released in December and backed by many intellectuals.

After signing, Professor Xia, the economist, said he was forced to resign from positions at two
research institutes. His fellow signer, He Weifang, a celebrated law professor, was transferred to
an obscure college in China's far west. Professor He's exile was news overseas. But much like the
coming anniversary of the Tiananmen Square crackdown, it drew little notice from students.

One student defended the professor with an anonymous post on the campus's computer bulletin board.
``The day will come,'' he wrote, ``when Professor He can go where he wants.''

\section{Climate Bill Clears Hurdle, but Others Remain}

WASHINGTON -- The House Energy and Commerce Committee, splitting largely along party lines, approved
on Thursday the most ambitious energy and global warming legislation ever debated in Congress.

The bill's passage, on a 33-to-25 vote, served as a bookend to a week that began with President
Obama's announcing a deal with auto manufacturers to impose tough new mileage and emissions
standards for all cars and trucks sold in the United States starting in 2012.

With progress on this legislation and his own executive actions, Mr.~Obama is assembling the pieces
of a credible American package to take to Copenhagen later this year as United Nations negotiators
gather for talks on a new global warming treaty.

But the energy legislation passed on Thursday still faces a tortuous path through several more House
committees before it can be brought up for a vote later this year. In the Senate, leaders say they
lack the votes to pass the bill as it is now written.

Representatives Henry A.~Waxman of California and Edward J.~Markey of Massachusetts, both Democrats,
sponsored the bill. Mr.~Waxman called it ``decisive and historic.''

``When this bill is enacted,'' Mr.~Waxman said, ``we'll break our dependence on foreign oil, make
our nation a leader in clean-energy jobs and technology and cut global warming pollution.''

The objections of the Republican opponents were summed up in the words of Representative Mike Rogers
of Michigan, who said the bill would mean sharp increases in energy costs and the loss of millions
of jobs.

``This is the biggest energy tax in the history of the United States,'' Mr.~Rogers said.

Mr.~Obama did not play a major public role in the committee's work, but intervened quietly on
several occasions, calling nervous conservative Democrats to assure them that a vote for the bill
would not hurt them politically. Two weeks ago, he gathered all of the panel's Democrats at the
White House to urge them to set aside their differences to produce a bill that met his goals of
energy conservation and global warming abatement.

The measure approved by the House committee runs more than 930 pages. It establishes a cap-and-trade
program to control climate-altering emissions; dictates an increase in the use of renewable energy
sources; and sets new efficiency standards for buildings, lighting and industrial facilities. It
calls for a 17 percent reduction in emissions of heat-trapping gases from 2005 levels by 2020 and 83
percent by 2050.

It was the object of one of the biggest lobbying campaigns of any piece of environmental
legislation, with millions of dollars spent on both sides in the months leading up to Thursday's
vote. Lawmakers heard from former Vice President Al Gore, local utility companies, hunters and
fishermen, national environmental groups, agricultural interests and the coal, oil and natural gas
industries.

Much more is to come as the bill, known as the American Clean Energy and Security Act, is taken up
by the Ways and Means Committee and, possibly, House panels on agriculture, transportation and
natural resources.

The measure has critics and supporters on the left and right. Many large manufacturers and the
majority of the big utilities support it, for now, as the least-bad option for what is considered
the inevitable imposition of controls on heat-trapping gases like carbon dioxide.

The National Association of Manufacturers and the United States Chamber of Commerce opposed it,
although some members of both groups broke with leadership to endorse the measure. These two big
lobbies argued that it was a burdensome tax on business that would drive companies and jobs overseas
while doing little to address climate change.

The major environmental advocacy groups are, for the most part, enthusiastic backers. ``Clearly,
it's an unprecedented step forward in the critical effort to create millions of clean-energy jobs,
reduce our dependence on oil and protect our planet from global warming pollution,'' said Gene
Karpinski, president of the League of Conservation Voters.

But some groups, like Greenpeace, Friends of the Earth and Public Citizen, expressed concern that
the bill as drafted gave away too much to polluting industries and depended excessively on
hypothetical reductions in emissions of heat-trapping gases from developing countries.

``Despite the best efforts of Chairman Waxman, the decision-making process was co-opted by oil and
coal lobbyists determined to sustain our addiction to dirty fossil fuels,'' said a joint statement
from the groups. ``The resulting bill reflects the triumph of politics over science, and the triumph
of industry influence over public interest.''

When Mr.~Waxman first unveiled his plan in late March, at least a dozen of the panel's 36 Democrats
had qualms of it. These so-called Brown Dogs were mainly from states dependent on coal for power and
manufacturing for jobs, and needed assurance that their constituents would be protected.

In weeks of closed-door negotiations with these Democrats, Mr.~Waxman doled out billions of dollars
worth of free pollution permits, known as allowances, to cushion any price shock caused by imposing
a cap on emissions of heat-trapping gases.

In the end, 85 percent of all pollution allowances were given at no cost for various purposes,
including compensating energy-intensive industries, state governments, oil refiners and low-income
households, at least in the early years of the program. Mr.~Obama's position during the presidential
campaign was that all of the permits should be auctioned, not given away, but the White House did
not object to Mr.~Waxman's generous allocations.


\section{Potential Risk in a Move to Give Investors More Say on Boards}

The Securities and Exchange Commission is worried about shareholder democracy.

Again.

The commission is asking for comment on proposed rules that would make it easier for institutional
shareholders to propose board candidates to be listed on the proxies that every public company sends
to its shareholders.

If a rule is adopted -- and it looks as if Mary L.~Schapiro, the chairwoman, and her two Democratic
colleagues are determined to make changes before next spring's proxy season -- then shareholders
will get something they have never seen before: ballots with more candidates than jobs to be filled.

This is an issue that has been around for decades, and that may have cost one of Ms.~Schapiro's
predecessors his job. Corporate managements hate the idea, but some institutional investors,
including public pension funds, love it, claiming they will be able to improve corporate governance.
Companies warn the rule will lead qualified director candidates to refuse to serve rather than face
the possibility of rejection by the shareholders.

While Washington has debated the idea off and on since at least 1977, the world has changed. On the
one hand, the need for such a rule has ebbed as markets and technologies made it a lot easier and
cheaper for dissatisfied investors to try to replace individual directors or whole boards by running
traditional proxy fights.

On the other hand, the expansion of derivatives markets has made it easy to separate a share's vote
from the economic interest that used to be a fundamental part of owning a share. There have been
shareholder votes dominated by ``owners'' whose own financial interests ran counter to those of
shareholders. Now there are worries that bondholders, whose votes can be critical in bankruptcy
reorganizations, may want a company to fail so they can collect on their much larger positions in
credit-default swaps.

In retrospect, shareholder democracy issues should have been considered when these derivative
securities were being invented and allowed to trade. But they were not.

Just why that happened is a question that Congress and the Obama administration should ponder as
they revamp the American regulatory structure. Perhaps the S.E.C. would have considered the effects
of divorcing votes from economic interest when the derivatives began trading. But some of these
products were completely unregulated, and others went through the Commodities Futures Trading
Commission, which in recent years was uninterested in regulating anything.

In other words, if they are going to make it easier for shareholders to cast meaningful votes,
perhaps they should also do something to assure that the voters have no financial interest in the
company's failure. Henry T.~C. Hu, a law professor at the University of Texas who has written
extensively on this question, told me this week that it might make sense to require those casting
large blocks of votes to certify they did not have an economic interest contrary to that of the
company. But defining that could be treacherous.

The proposal put forth by the S.E.C. this week would let any holder, or group of holders, of at
least 1 percent of the shares in large companies nominate a director candidate. To do that, the
nominating group would have to certify that it was not interested in a takeover of the company, and
that it had owned the shares for at least one year.

How many directors could be so nominated? Up to a quarter of the board, so two seats could be
contested if the board had eight members. And if two groups want to make nominations for one seat?
The first one in gets on the ballot.

That holding period could raise its own questions. If an institution lends out shares for shorting,
it at least theoretically also lends out the voting rights. Does that loan stop the holding period?
If so, would that cause some institutions to stop lending shares, making it harder to short stocks
and perhaps contributing to market volatility?

If a director contest does emerge, the proxy sent out by a company might end up listing five
candidates for four seats. We're all used to that in political campaigns, but it is unheard of in
corporate contests. In conventional proxy fights, each side sends out its own ballot. If you support
management, sign the white card, which lists only its candidates. If you support the insurgents,
sign the blue card.

Decades ago, it was very expensive to mount such a proxy fight, given the cost of mailings to
individual shareholders who collectively dominated many companies. Now, however, a majority of votes
is often held by a relatively small number of institutions, and a proxy fight can be mounted for
much less, with communications done by phone and e-mail with major holders. The individual investors
may never be contacted.

A few years ago, it appeared to be a sure thing that a much milder rule on shareholder access would
be adopted. That rule, endorsed by William H.~Donaldson, then the S.E.C. chairman, set up a lot of
hoops to be leaped through before an institutional shareholder could nominate a rival candidate to
appear on the company's proxy.

Mr.~Donaldson, a Republican, had the support of two of the other four commissioners and could have
forced through a rule. But his supporters were Democrats. Republicans bemoaned his lack of party
loyalty and he did not bring the issue to a vote. He was eventually succeeded by Christopher Cox,
who pronounced ``consensus'' to be his goal. Shareholder access was dead.

Since then, some companies have made it easier for shareholders to reject directors by requiring
majority voting for elections. If an unopposed candidate is opposed by more owners of most shares
being voted, no one is elected to the seat. A pending rule of the New York Stock Exchange would bar
incumbents from being backed by brokers who are voting shares owned by clients who did not bother to
vote. Had that rule been in effect this year, two Citigroup directors would have been denied
re-election.

In principle, shareholder access seems to be an obvious good idea. Companies do, after all, belong
to their shareholders. We will see if this does become, as its advocates forecast, a rarely used
instrument to change companies whose boards have failed to perform adequately, or if it will lead to
a series of bitter battles -- sometimes instituted by institutions with ulterior motives -- that result in divided boards and divert both directors and managers from running companies to benefit all shareholders.

\section{`Idol' Strives to Sustain Its High Note}

LOS ANGELES -- The confetti had not yet stopped falling in the Nokia Theater here on Wednesday night
when the conspiracy theories began to emerge: The producers of ``American Idol'' wanted Kris Allen
to win. The vote was rigged. Fox was determined to have a fresh-faced, heartland-appropriate,
eyelinerless Idol as the public face of the most popular series on television.

So too did the requiem begin for ``American Idol.'' The ratings were the lowest in the show's
history for a finale, the headlines read. The experiment of adding a fourth judge was a failure,
critics proclaimed. Simon Cowell may be leaving the show, spelling its doom.

Don't start digging a grave just yet however. Those who think ``American Idol'' is losing its grip
on the top of the ratings are probably fooling themselves, at least according to Fox.

``It's ridiculous how big this show is,'' Mike Darnell, president for alternative programming at Fox
Broadcasting, said in an interview on Friday. ``It has had just as big an impact as ever this
season, maybe bigger in terms of the amount of cultural talk it has generated.''

The series certainly seemed to be on the tip of tongues around the country this season, with a field
of Top 5 performers that was arguably the show's strongest ever. Adam Lambert, in particular, drew
the most attention, dividing viewers while also drawing admiration for the energy he injected into
performances.

Mr.~Allen, meanwhile, played the part of Mine That Bird, the unsung colt who climbed unnoticed from
the back of the pack to Kentucky Derby victory.

With even the ``Idol'' judges seeming to expect Mr.~Lambert to win, it did not take long for the
hothouse of the blogosphere to generate numerous theories about the forces that had aligned against
him.

The show's host, Ryan Seacrest, it was noted, mentioned multiple times during the finale that nearly
100 million votes were cast, but he gave no hint about the margin separating the two finalists. Last
year he clearly stated that the winner, David Cook, received 56 percent of the 97.5 million votes
cast, a 12 million-vote margin of victory. And earlier this season Mr.~Seacrest gave details about
the vote margin between the bottom two contestants.

Many Web commentators seized on questions about Mr.~Lambert's sexuality to claim that a homophobic
campaign was being waged against him. (Mr.~Lambert has not stated his sexual orientation.) But that
argument ignores Mr.~Lambert's rise to the final round. Presumably someone intent on eliminating him
for reasons not related to his singing could have done so earlier in the season.

The Fox affiliate in Little Rock, Ark., reported on Wednesday that ``an extraordinary number of
votes'' -- 38 million in some reports -- came from Arkansas, a state with a population of less than
3 million. That detail quickly made its way around the Internet. The station later retracted that
news, saying that it was ``based on incorrect information provided to one of our reporters.''

Mr.~Seacrest also told the audience on Tuesday that the voting would continue for ``at least four
hours'' after the end of the show, an open-ended deadline that raised as many questions as it
answered. Would the lines stay open only until the ``right'' candidate was ahead?

Mr.~Darnell said that the results were not manipulated, and that Fox and the producers had no
preferences regarding the winner. ``There are no conspiracies, and there is no plot,'' he said.
``We're happy to have America tell us who the winner is. And they have always gotten it right.''

What Mr.~Allen had going for him, most likely, is that he was from the South. Consider the eight
winners of ``American Idol'' and their home states: Kelly Clarkson, Texas; Ruben Studdard, Alabama;
Fantasia Barrino, North Carolina; Carrie Underwood, Oklahoma; Taylor Hicks, Alabama; Jordin Sparks,
Arizona; David Cook, Missouri; Kris Allen, Arkansas.

It is difficult to tell what part audience demographics played in Mr.~Allen's victory. Many
commentators claimed that texting teenagers favored him, but if that was true, they probably also
would have favored David Archuleta over Mr.~Cook last season.

And the ``American Idol'' audience is getting older. In 2004 viewers 6 to 24 made up 24 percent of
the ``Idol'' audience, according to Nielsen Media Research, while viewers over 55 accounted for 18
percent. By this year those numbers had flipped: the 6-to-24 category measured only 17 percent of
the total audience, while the over-55 contingent was 27 percent of the total, according to Nielsen.

Fox noted that although the Nielsen ratings for the finale were the lowest in the show's eight years
among 18-to-49-year-olds, the total number of viewers for the finale -- 28.8 million -- was the same
as the third-season finale, when Ms.~Barrino outdueled Diana DeGarmo.

``You can't compare this year's ratings to anything on television seven years ago,'' Mr.~Darnell
said, because everything on television is drawing fewer viewers. The spread between ``Idol'' and the
No.~2-rated show has only widened in recent years.

As for Mr.~Cowell, he is under contract for next season, as is Randy Jackson, but Paula Abdul and
Kara DioGuardi are not. Mr.~Darnell would not say whether the show would return to three judges or
keep four.

``No decision has been made,'' he said. ``We made an enormous amount of changes this year. Now we'll
go back and consider which things worked and which ones didn't.''

Or, as Mr.~Seacrest would say: The answer is -- after the break.

\section{Payoff Over a Web Sensation Is Elusive}

Susan Boyle, the frumpy Scotswoman who became a worldwide singing sensation last month, may wind up
as the winner this week of ``Britain's Got Talent,'' the hit ITV show.

After a six-week absence, she returned on Sunday night to sing ``Memory'' from the musical ``Cats,''
wowing the crowd and advancing to Saturday's finale. The producers immediately posted her
performance on the Internet for the rest of the world to see.

She has already won a popularity contest on YouTube, where videos of her performances in April have
been viewed an astounding 220 million times.

But until now, her runaway Web success has made little money for the program's producers or
distributors.

FremantleMedia Enterprises, a production company that owns the international digital rights to the
talent show, hastily uploaded video clips to YouTube in the wake of Ms.~Boyle's debut, but the clips
do not appear to be generating any advertising revenue for the company. The most popular videos of
Ms.~Boyle were not the official versions but rather copies of the TV show posted by individual
users.

The case reflects the inability of big media companies to maximize profit from supersize Internet
audiences that seem to come from nowhere. In essence, the complexities of TV production are curbing
the Web possibilities. ``Britain's Got Talent'' is produced jointly by three companies and
distributed in Britain by a fourth, ITV, making it difficult to ascertain which of the companies can
claim a video as its own.

Before the current season of the talent show started on April 11, the parties tried to cut a
distribution deal with YouTube, but they could not agree on terms, according to two people with
knowledge of the talks. The people asked for anonymity before they would discuss confidential
negotiations.

YouTube, a unit of Google, has been keen to make money from its hulking library of online video by
signing contracts with copyright owners and sharing the revenue from ads it sells before, during,
after and alongside the videos. Major media companies have shown varying degrees of interest in
these deals, in part because they are reticent to split much money with Google.

Then Simon Cowell, an ``American Idol'' judge who is also a producer and a host of ``Britain's Got
Talent,'' helped introduce Ms.~Boyle to the world.

Her performance was a made-for-TV fairy tale: a dowdy 48-year-old makes awkward jokes, the audience
engages in a collective eye-roll, then the performer shocks everyone by bursting into a soulful,
Broadway-worthy rendition of ``I Dreamed a Dream.''

Cut to the amazed faces in the theater, hear the judge Piers Morgan call her singing ``without a
doubt the biggest surprise I have had in three years on this show,'' and cut to commercial.

On YouTube, though, where the segment was viewed by more people than could ever have witnessed it on
TV in Britain, there were no commercials. The tens of millions of views swiftly brought YouTube and
the producers back to the negotiating table, according to the people with knowledge of the talks,
and soon they reached a deal for video clips.

YouTube was especially interested in a deal, according to the people with knowledge of the talks,
because the company was essentially losing money by serving every video stream without recouping any
of the costs.

FremantleMedia, which had registered YouTube accounts for the next several seasons of ``Britain's
Got Talent'' in advance, uploaded dozens of clips from the show in late April. But American viewers
are not seeing ads on the video pages, suggesting that the companies still do not see eye to eye.

FremantleMedia ``is investigating the best routes to monetize the channel in conjunction with
relevant partners,'' said a spokeswoman, Belinda Thomas, who said the company would not comment
further.

The production companies and YouTube worked through the weekend on a more comprehensive deal, one of
the people with knowledge of the talks said. The deal would enable FremantleMedia to place ads
against unofficial copies of the show, using YouTube's ``Content ID'' system, which companies like
Universal Music already use. For now, the copies simply show a message directing users to the
official talent show channel managed by FremantleMedia.

``We're glad to be helping Britain share its talents with the rest of the world,'' a YouTube
spokesman, Ricardo Reyes, said. ``It's up to our partners to decide what to do with their videos on
YouTube.''

How much money have the parties lost? In the days after Ms.~Boyle's debut, The Times of London
published what it called a ``crude estimate'' suggesting that the parties involved had left \$1.87
million on the table.

That is based on 75 million streams of the various clips of Ms.~Boyle, which the newspaper estimated
could get \$20 to \$35 for every 1,000 views in the United States, and more than that in Britain.

While other TV networks act quickly to remove videos when users upload them without copyright
permissions, ITV has ``nonexistent piracy enforcement on YouTube,'' said David Burch, a marketing
manager at TubeMogul, an online measurement firm.

The broadcaster and producers allowed the copies to stay online because they created buzz for the
program. The clips have received more than a half-million user comments.

The view counts continued to grow as people awaited Ms.~Boyle's next performance. Visible Measures,
a company that tracks online video placements, said Ms.~Boyle was responsible for the
fastest-growing viral video in the roughly five-year history of Web video. Only three other videos
have received more clicks, said the company, which tracks viewing across about 150 sites. (YouTube
is the biggest by far.)

Matt Cutler, the vice president for marketing and analytics at Visible Measures, said the level of
interest was ``off the charts.''

``On TV, watching the content is the end of the experience. Online, watching the content is the
beginning of the experience,'' Mr.~Cutler said.

The history of viral videos has shown that when new clips about a subject become available -- in
Ms.~Boyle's case, her new performance on Sunday -- it ``actually boosts the viewership of the
existing assets,'' Mr.~Cutler said.

Six hours after the new performance, dozens of copies were already circulating on YouTube.

\section{AT\&T May Have Swayed 'Idol' Results}

LOS ANGELES -- AT\&T, one of the biggest corporate sponsors of ``American Idol,'' might have
influenced the outcome of this year's competition by providing free text-messaging services and
lessons in casting blocks of votes at parties organized by fans of Kris Allen, the Arkansas singer
who was named the winner of the show last week.

Representatives of AT\&T, whose mobile phone network is the only one that can be used to cast
``American Idol'' votes via text message, provided the free text-messaging services at two parties
in Arkansas after the final performance episode of ``American Idol'' last week, according to the
company and people at the events.

There appear to have been no similar efforts to provide free texting services to supporters of Adam
Lambert, who finished as the runner-up to Mr.~Allen.

Since then, angry supporters of Mr.~Lambert have flooded online chat boards with messages claiming
irregularities in the competition's voting procedures.

Officials of AT\&T and Fox Broadcasting declined to discuss the situation. Details of the voting
support were first reported last week in an article in the Arkansas Democrat-Gazette.

In addition to the free texting services, representatives of AT\&T also helped fans of Mr.~Allen at
the two Arkansas events with instructions on how to send multiple, simultaneous ``power text''
messages into the ``American Idol'' voting lines.

By sending 10 or more text messages at the press of a single button, ``power texts'' have an
exponentially greater effect on voting than do single text messages or calls to the show's toll-free
phone lines.

The efforts appear to run afoul of ``American Idol'' voting rules in two ways. The show broadcasts
an on-screen statement at the end of each episode warning that blocks of votes cast using
``technical enhancements'' that unfairly influence the outcome of voting can be thrown out.

And the show regularly states that text voting is open only to AT\&T subscribers and is subject to
normal rates.

In a statement issued Tuesday, a spokesman for AT\&T said: ``In Arkansas, we were invited to attend
the local watch parties organized by the community. A few local employees brought a small number of
demo phones with them and provided texting tutorials to those who were interested.''

\section{Game Maker Trims Revenue Forecast}

The video game publisher Take-Two Interactive reported stronger-than-expected quarterly results on
Tuesday but cut its revenue forecast for the year and delayed the release of two games.

The company said it would delay the release of Red Dead Redemption and Mafia II from the fourth
quarter to the first half of fiscal 2010 and would look for ways to cut costs.

Take-Two said its fiscal year performance would be strongly weighted to the fourth quarter, after
the release of Grand Theft Auto: Episodes From Liberty City on the Xbox 360.

The company said lost \$10.1 million, or 13 cents a share, in the fiscal second quarter ended April
30, compared with a profit of \$98.2 million, or \$1.29 a share, in the year-ago period.

Excluding certain items, Take-Two lost 4 cents a share, better than analysts' average estimate for a
loss of 14 cents a share, according to Reuters Estimates.

Revenue tumbled 57 percent to \$229.7 million, compared with the Wall Street forecast of \$216.6
million. In the second quarter of last year, the company had released its wildly successful Grand
Theft Auto IV game.

The company reaffirmed its forecast for the fiscal year ending Oct.~31 of a profit, excluding items,
of zero to 20 cents a share. But it lowered its revenue forecast to a range of \$1.05 billion to
\$1.15 billion. It had previously estimated revenue of \$1.1 billion to \$1.25 billion.

\section{In Sports Business, Too Many Hopefuls for Too Few Positions}

Jason Martin entered the University of South Carolina's sport and entertainment management program
hoping to become the next Jerry Maguire, the mythical super agent. Instead, like most people trying
to break into the sports industry, he is looking at years of low-paying, unglamorous jobs like
selling ads, tickets and hot dogs.

That, of course, assumes he can find a job. In January, he started a four-month unpaid internship at
the Family Circle Cup, a women's tennis tournament, helping sponsors place their ads. He hoped his
internship would lead to a salaried job.

Not this year. Because of the recession, Martin, who graduated this month with thousands of other
sports management majors across the country, is continuing his job search at home in Virginia. If no
full-time work turns up by September, he will enroll in a master's degree program.

``Graduate school is definitely Plan B,'' said Martin, 21, who is working this summer as a swimming
instructor. ``My original intention was to get a job, but with the economy, there's so many people
who just graduated who can't even get a cup of coffee with a prospective employer.''

For decades, the sports industry has been largely impervious to the economic cycle. Through booms
and busts, leagues and tournaments expanded, new stadiums were built and attendance and television
viewership hit records. Revenue from suite sales, naming rights and television contracts boomed.

But Martin and other graduates are finding out that the industry's growth is slowing, if not
reversing. Students are receiving fewer job offers this spring or are being forced to take
internships in lieu of a salaried position. Many of those internships are unpaid. The worry, their
professors say, is that austerity may become the norm, forcing students to scale down or abandon
their ESPN-fueled dreams.

``I used to teach that sport was recession proof, but this recession proved me wrong,'' said Gary
Sailes, who runs the undergraduate program in sports marketing and management at Indiana. ``I tell
students that this is a good time to stay in school.''

The number of jobs related to spectator sports has risen steadily during the last few decades.
According to figures from the Bureau of Labor Statistics, 138,700 people work in the spectator
sports industry, 9.8 percent more than in 2002.

There is no accurate tally of how many jobs have been lost in the last year or so, but the anecdotal
evidence is grim. The N.F.L. has cut nearly 200 jobs. The N.B.A. has eliminated 10 percent of its
staff. And the United States Olympic Committee laid off more than 50 workers.

The L.P.G.A. has dropped several tournaments, and Honda ended its Formula One sponsorship. The Jets
will furlough some employees for two weeks, and the Cleveland Browns, the Denver Broncos and the
Washington Redskins have cut jobs.

Many other teams and tournaments, stung by declines in attendance and sponsorship dollars, have
stopped hiring. ESPN will not fill 200 vacant jobs. The Arena Football League canceled its season.

This is bad news for the 300 or so universities that offer sport management degrees. Every year,
they churn out thousands of graduates who, even in good times, are willing to work for low pay in
return for the chance to work around athletes and arenas.

The teams, leagues and others in the sports industry have taken advantage of their willingness to
make financial sacrifices, and may continue to do so.

``Sport management is one of the sexy industries, and it's very hard to discourage students from
joining the industry,'' said Mark McDonald, the director of internships at the University of
Massachusetts's sport management program. ``I feel for the students who haven't done the
preparation. They are in for a rude surprise.''

Going to college to prepare for a job in sports is relatively new. Years ago, many jobs were filled
though word of mouth or serendipitously. Marty Appel, the Yankees' assistant public-relations
director in the early 1970s, wrote to the team in 1967 asking for a summer job. He was lucky;
someone was needed to answer Mickey Mantle's mail.

``It was very much like a fly-by-the-seat-of-your-pants operation,'' said Appel, who figured there
were only 40 people working in the team's front office then.

By the late 1970s, Appel interviewed graduates of some of the first sport management programs who,
he said, were better prepared to handle the growing list of demands on ball clubs.

Walter O'Malley, who owned the Dodgers in Brooklyn and Los Angeles and who foresaw the need for
better trained employees, helped persuade Ohio University to start the first degree-granting sport
management program in 1966.

The number of programs jumped fourfold in the 1980s as the industry grew along with ESPN and Title
IX, according to Jacquelyn Cuneen, who teaches sport management at Bowling Green. Dozens of
universities, looking to boost enrollment, rebranded their physical education departments.

Sport management programs have also been broadened to teach fitness, tourism, recreation and
hospitality, prompting critics to claim that many programs do not provide students with the skills
to succeed at anything other than the lowest-level jobs.

More worrisome is the realization that even at rigorous programs, students are getting about half as
many job offers this year compared with other years, according to Dallas Branch, who teaches sport
management at West Virginia. Yeoman's work is demanded of graduates even from the most prestigious
universities like Columbia, which started a master's level program for midcareer students in 2006.

While attending classes there, Sean Mysel is working this summer as a stadium manager for the Sussex
Skyhawks, a minor league baseball team in New Jersey. But he has spent about \$25,000 on tuition and
\$75,000 on living expenses to attend Columbia, which he hopes will help him reach his goal of
running a major league ballpark.

``I know the job market is difficult for a while, but I don't feel I have to hit the panic button,''
said Mysel, 31. ``At a minimum, it will take 10 or 15 years.''

Graduate students, though, are in a better position than undergraduates because they often have
several years of work experience. Kati Karottki, who graduated last weekend from Massachusetts with
a bachelor's degree in sport management, did well in college and worked as a research analyst at a
company that measures sponsorships.

But as her senior year began, she planned to single out several companies as potential employers.
Then Lehman Brothers collapsed in September.

``I really had to redo everything,'' she said.

She reached out to alumni and had half a dozen interviews. But over and over, she heard the same
thing: there are other candidates with master's degrees with more experience who are willing to work
for little. This week, she is traveling to Bristol, Conn, to interview at ESPN for an analyst's job
focused on audience research. She is not getting ahead of herself.

``They have so many applicants and people to choose from, just getting to an in-person interview,
that's a big achievement,'' she said. ``My classmates have resigned themselves to doing another
internship or taking time off to travel.''

\section{Chinese President Meets Leader of Taiwanese Party}

BEIJING -- President Hu Jintao of China hosted the leader of Taiwan's governing party in Beijing on
Tuesday, as both sides promised closer economic ties and a continued thaw in relations between the
rivals.

China promoted the visit as a step toward more investment in Taiwan's technology sector, while
Taiwanese officials said they hoped to benefit from an influx of mainland Chinese tourists.

The two sides have not begun to address the political future of Taiwan, a self-governing island that
mainland China claims as sovereign territory. China maintains missile deployments that Taiwan says
are aimed at bullying it into compliance with Beijing's wishes.

But since the Kuomintang's victory in presidential elections last year, replacing the
independence-leaning Democratic Progressive Party, Taiwanese and mainland officials have worked to
show some concrete results.

``The Chinese haven't pulled any missiles out. It's not a peace agreement, but a tempering down'' of
old hostilities, Bruce Jacobs, who directs Taiwan research at Australia's Monash University, said in
an interview.

That said, he added, ``Things are going relatively quietly.''

Ma Ying-jeou, who became president last May, still rejects Beijing's claims of sovereignty. But he
has worked for closer economic ties and regular dialogue between party officials in both places.

The stop in Beijing on Tuesday by Wu Poh-hsiung, the Kuomintang Party leader, was not by itself
groundbreaking. He has made others in recent years. But other events underscore the extent to which
the political dynamic is changing.

Last week Chen Chu, the highest-ranking elected official of the Democratic Progressive Party, or
D.P.P., visited Beijing. Ms.~Chen, the mayor of Kaohsiung, Taiwan's second largest city, stressed
that she was visiting the mainland to promote her city, not to signal any political retreat from the
party's independence stand. Trips to the mainland by D.P.P. officials are nevertheless rare, and
Mr.~Ma quickly cited it as a sign that cross-straits relations were improving.

Beijing has also shown some diplomatic flexibility. In April, the Chinese government allowed Taiwan
to participate as an observer at the World Health Assembly, the governing body of the World Health
Organization.

It was the first time that Taiwan had been granted observer status at a United Nations body since it
lost the China seat in 1971, albeit not formally as Taiwan, but as Chinese Taipei.

The two sides have reached some concrete agreements as well. This year, the first direct air service
between Taiwan and mainland China was started, scrapping the traditional subterfuge of routing all
traffic through Hong Kong. More recently, Taiwan increased the daily quota of visitors from China to
3,000, a tenfold increase.

The governments are also discussing increased cooperation on law enforcement issues like joint
criminal investigations and extradition.

The meeting Tuesday with Mr.~Hu was devoted to economic issues. Taiwan has announced plans to allow
Chinese businesses and banks greater leeway to invest in the island.

But critics have questioned whether the Chinese should be allowed to own stakes in strategic assets
like airports and harbors, and some have charged that Mr.~Ma is clearing the way for a Chinese
takeover of Taiwan's economy.

\section{Modifying Conciliatory Stance, South Korea Pushes Back Against the North}

SEOUL, South Korea -- After years of incendiary threats and weapons tests, the world has become
accustomed to provocations from North Korea. More unusual this time is that South Korea is pushing
back.

After North Korea's nuclear test last week, South Korea swiftly moved to join an American-led effort
to search ships for unconventional weapons, despite threats of war from the North if its ships were
boarded. South Korea has also lobbied China and Russia to take a unified stance against the North,
and supported calls at the United Nations Security Council for new sanctions.

The assertive moves by South Korea's right-leaning and generally pro-American president, Lee
Myung-bak, are a sharp departure from the conciliatory stance that the South took in previous
standoffs, particularly after the North's first nuclear, test three years ago, when it continued
economic aid and investment. South Korea's reluctance to challenge threats from its neighbor, along
with China's more accommodating policies, made it difficult in the past to enact or enforce tough
sanctions on North Korea. Now, the South's decision to stop acting as the North's leading protector
may open a wider range of options in responding to the nuclear and missile tests, analysts say.

``South Korean aid gave North Korea leverage against the United States,'' said Choi Jin-wook, a
North Korea specialist at the Korea Institute for National Unification, a state-run research
organization. ``Lee Myung-bak has emptied their pockets and weakened their bargaining position.''

It remains to be seen how much that will translate into punitive actions against the North at the
United Nations, where the five permanent Security Council members -- Russia, China, the United
States, Britain and France -- have been joined by Japan and South Korea in trying to hammer out a
new resolution that includes sanctions.

There is broad agreement on the need to send a tough message to North Korea, senior diplomats have
said, but less consensus on what that should be. Given the sensitivities of China, which is the
North's one ally, the resolution will probably not involve any new sanctions. Instead, it may impose
some of those approved in the resolution after the 2006 nuclear test but never enacted because the
North cooperated in entering six-party talks on its weapons program.

That earlier resolution focused on cutting trade and financing for North Korea's military and
hampering the ruling elite with a ban on the import of luxury goods. No one expects rapid action,
however. ``It is not a normal country, so we have to act cautiously,'' Zhang Yesui, the Chinese
ambassador to the United Nations, said in an interview on Sunday. ``It is a very closed place, so
they act abnormally.'' It is not even clear that China has the influence to sway the North Korean
government, he added.

It is also unclear how far South Korea is willing to push the North, whose troops and artillery sit
just 35 miles north of Seoul along the heavily fortified demilitarized zone. While Mr.~Lee did gain
office last year promising a tougher line, most South Koreans still abhor tensions with the North
and are fundamentally in favor of reconciliation, said North Korea experts.

``If North Korea actually tests South Korea's resolve, it will put South Korea in a very tough
spot,'' said Kim Il-young, a North Korea specialist at Seoul's Sungkyunkwan University. Indeed,
South Korean officials have been careful to cast their new approach not as a drastic departure but
as more a shift in tone. They say the South remains willing to offer economic aid to the North, but
only if it responds with good-faith efforts to curtail its weapons programs.

They also say that they are not trying to provoke the North, which has seemed particularly volatile
since its leader, Kim Jong-il, appears to have had a stroke in August. ``This is not at all a
hawkish or hard line,'' said Kim Tae-hyo, the secretary to Mr.~Lee for national security. ``We still
seek contacts with the North. Where we differ is that we demand mutual respect and reciprocity.''

Still, by moving the nation closer to Washington, the South Korean president's approach has deprived
North Korea of one of its usual policy options in such standoffs: driving a wedge between Seoul and
its traditional allies in Washington and Tokyo.

Emblematic of Seoul's new tilt toward the United States was its sudden decision on Tuesday to join
the Proliferation Security Initiative, the effort to stop global trafficking in unconventional
weapons, after years of refusing to do so for fear of riling the North. North Korea responded by
reiterating that it would consider any search of its ships to be an act of war.

Even during the current crisis, however, South Korean officials said they were still trying to reach
out to the North. They also said they were talking with China, Russia, Japan and the United States
about putting together a ``grand plan'' of aid and other economic incentives to persuade the North
to dismantle its nuclear program.

As a possible model, they cite Ukraine, which relinquished its Soviet-made nuclear arsenal in
exchange for a huge aid package and entree to the West. However, with the North's fear of the
instability that greater openness could bring, many analysts doubt that even a generous aid package
could coax it to end its weapons program.

The emerging view here is that last Monday's nuclear test and the ensuing series of test missile
launchings reflect North Korea's attempt to ensure its survival by creating an independent nuclear
deterrent. The displays of technological achievement are also seen as a bid to create domestic
support for the ill leader and his moves to install a successor, South Korean officials said.

South Korean officials play down their role in provoking the North, saying officials in Pyongyang
are engaging in the same brand of nuclear brinkmanship they have used in previous crises, though
with a markedly faster pace. They say that despite the North's recent threats against the South, the
main target audience remains the United States, with which the North wants to hold bilateral
negotiations on a new footing as a nuclear nation.

``You can't connect the dots between South Korean policy changes and the scale and pace of North
Korea's current provocations,'' said Wi Sung-lac, the South Korean Foreign Ministry's top negotiator
on Korean peninsula security affairs.

\section{Civil Liberties Within Limits After 12 Years of Beijing Rule}

HONG KONG -- It was a raucous display of free speech outside Hong Kong's Legislative Council last
week: construction workers demanding increased spending on public works, retirees agitating for
heftier pensions, and legislators, wearing black T-shirts printed with tanks, calling on the Beijing
government to apologize for the bloodshed in Tiananmen Square two decades ago. In the 12 years since
it passed from British to Chinese rule, Hong Kong has remained a bastion of civil liberties unknown
in mainland China, under an arrangement dubbed ``one country, two systems.''

The result has been the continuation of a freewheeling press, an independent judiciary and a
well-oiled bureaucracy. On Thursday, tens of thousands are expected to turn out for a candlelight
vigil in Victoria Park here to commemorate the 1989 Tiananmen crackdown, in which hundreds of
students advocating democracy were killed. In the rest of China, any mention of the events at
Tiananmen Square has been banned in the news media or public discourse.

But many democracy advocates and civil libertarians here are increasingly anxious about whether
laissez-faire Hong Kong can maintain its independence from Beijing's authoritarian grip and its
distinct identity as an amalgam of Western and Chinese sensibilities.

Last year, Beijing postponed direct elections -- to 2017 for the chief executive and 2020 for the
full legislature -- and its critics say China is wielding a heavier hand in Hong Kong's affairs.

A growing roster of overseas visitors whose politics irritate Beijing have been denied entry to Hong
Kong, and pro-China legislators have blocked efforts to include an uncensored account of Tiananmen
Square in high school textbooks.

Longtime advocates of democracy like Martin Lee warn that China is chipping away at Hong Kong's
autonomy by fiat or by co-opting business leaders and politicians.

On Saturday, Mr.~Lee, the founder of Hong Kong's Democratic Party, disclosed that he had been the
target of an assassination plot that he said the authorities foiled last August. He said the men
were arrested not long after he wrote an editorial accusing China of failing to live up to its
pledge to improve human rights.

``If you throw a frog into boiling water, it will jump out right away,'' Mr.~Lee said during an
interview in his office overlooking the High Court. ``But if you put the frog in warm water and cook
it slowly, it doesn't jump. We are being slowly cooked in Hong Kong, but hardly anyone is
noticing.''

Mr.~Lee and other democracy advocates have worried for years about Beijing's expanding influence
here. But in advance of the 20th anniversary of Tiananmen -- and of the July 1 anniversary of Hong
Kong's return to China in 1997 -- they are concerned about a new willingness by public officials to
openly back the mainland's view. They say that is a jarring development in a city where a million
people took to the streets in the summer of 1989 and where protests have been held every June 4
since then.

When asked last week if he supported exonerating the students who occupied Tiananmen Square, Hong
Kong's chief executive, Donald Tsang, told legislators that the episode was best forgotten. ``This
is something that happened a long time ago,'' he said. ``The national economy has grown and brought
prosperity to Hong Kong.'' He added that he thought his view ``represents the opinion of Hong Kong
people in general.''

Other local officials aligned with Beijing have gone further, claiming that no one died during the
crackdown, or that an armed response was warranted because student leaders were planning to kill
government soldiers. And this spring, Ayo Chan, the president of Hong Kong University's student
union, obliquely blamed the protesters for provoking the violence in Tiananmen.

Such statements do not go unchallenged, however. Angry students promptly voted Mr.~Chan out of
office, and Hong Kong's chief executive was forced to apologize. A poll by Hong Kong University last
week suggested that public sympathy for the Tiananmen protesters was high, with nearly 70 percent of
Hong Kong residents saying the Chinese government had erred in its handling of the demonstrations.

But those who closely watch the political culture here say reunification with China has begun to
slowly alter Hong Kong's unique ethos, even if the changes are hard to quantify and support for
democracy is still strong. Numbers tell part of the story: a decline in the number of Western
expatriates -- about 100,000 fewer since 1997 -- coinciding with a growing presence of mainlanders.

Last year, nearly 17 million mainland Chinese visited Hong Kong, compared with just over 2 million
in 1997. Shifting demographics have had an even greater impact on local universities. More than half
the postgraduate students studying here are from the mainland, up from barely one-third in 2003.

Like many other entrepreneurs here, Ronnie Chan, a billionaire whose company, Hang Lung Properties,
has expanded into the mainland, argues that Hong Kong can flourish only through closer ties to
China. In an interview last week, he said Hong Kong was far freer today than it ever was under the
British. If anything, he said, society could use a bit more restraint, especially when it comes to
the media. ``People were afraid the media would be curbed, but it's gone wild and become
irresponsible,'' said Mr.~Chan, 59, who attended college and graduate school in the United States.

Groups like the Hong Kong Journalists Association take a different view, saying the number of media
outlets willing to take on topics that might anger Beijing has been shrinking. Mak Yin-ting, a
freelance journalist and former chairwoman of the association, said the owners of more than half of
Hong Kong's media outlets -- many of whose owners have business interests in China -- have been
given advisory posts to the National People's Congress and People's Political Consultative
Conference, which rubber-stamp decisions made by the Communist Party.

The result, Ms.~Mak said, is that some reporters engage in self-censorship, while editors sometimes
bury stories that might be unflattering to Beijing. ``When your boss is a delegate to the National
People's Congress,'' she said, ``then you know it's better not to criticize China too loudly.''

\section{After Many Stumbles, the Fall of an American Giant}

It is a company that helped lift hundreds of thousands of American workers into the middle class. It
transformed Detroit into the Silicon Valley of its day, a symbol of America's talent for innovation.
It built celebrated cars, like Cadillacs, that became synonymous with luxury.

And now it is filing for bankruptcy, something that would have been unfathomable even a few years
ago, much less decades ago, when it was a dominant force in the American economy.

Rarely has a company fallen so far and so fast as General Motors. And while its bankruptcy appeared
increasingly likely in recent weeks, the arrival of the moment is still a staggering blow,
particularly for anyone with ties to the company.

``I never ever could have believed that one day this thing would go that way,'' said Jim Wangers, a
retired G.M. executive who was part of the team that developed the Pontiac GTO, and the author of
``Glory Days,'' about Pontiac's heyday in the muscle-car era of the 1960s. ``We were so
successful,'' he added.

Founded in 1908, G.M. ruled the car industry for more than half a century, with a broad range of
vehicles, reflecting the company's promise to offer ``a car for every purse and purpose.''

The expression ``What's good for General Motors is good for the country'' entered the lexicon, even
though it was a slight misquotation of Charles E.~Wilson, G.M.'s president in the early 1950s.

But then G.M. began a long and slow process of undermining itself. Its strengths, like the rigid
structure that provided discipline early on, became weaknesses, and it lost its feel for reading the
American car market it helped create, as Japanese automakers lured away even its most loyal buyers.

Only eight months ago, Rick Wagoner, then its chief executive, stood before hundreds of G.M.
employees to celebrate the company's 100th anniversary. ``We're a company that's ready to lead for
100 years to come,'' Mr.~Wagoner said.

Instead of leading, G.M. will instead be following other failed companies on a well-worn path into
bankruptcy court.

The moment will reverberate beyond G.M.'s epicenter in Detroit, to factory towns in other parts of
Michigan and in states like Indiana, Tennessee and Louisiana. It will even be felt on Fifth Avenue
in New York, where it built its financial headquarters, and Epcot at Walt Disney World in Florida,
where G.M. sponsors the Test Track Pavilion, a showcase of its latest cars.

G.M. factories churned out family cars, pickup trucks and memorable muscle cars with taut,
sculptured body panels that were rolling displays of American DNA.

A G.M. plant was a ticket to prosperity for the communities lucky enough to land one. G.M. literally
put Spring Hill, Tenn., on the map when it picked the town outside Nashville for its Saturn plant in
1985, prompting the hamlet to swell with new homes, motels and restaurants.

Now city officials around the country, including those in Spring Hill, nervously await phone calls
on Monday to tell them if their plants will be among the 14 G.M. is expected to announce it will
close in the latest round of cuts.

But even after its deep cuts, G.M. can still claim to be the country's largest automaker.

For G.M., that simple fact -- its sheer size -- was long used as a trump card to end debates. If the
critics were so right about all that was wrong with G.M., why did so many people buy its cars?

The company did have vast numbers of loyal buyers, but G.M. lost them through a series of strategic
and cultural missteps starting in the 1960s.

It bungled efforts in the 1980s to cut costs by sharing the underpinnings of its cars across
different brands, blurring their distinctiveness.

G.M. gave in to union demands in 1990 and created a program that paid workers even when plants were
not running, forcing it to build cars and trucks it could not sell without big incentives.

Its finance staff argued with product developers and marketers who pushed for aggressive spending on
new cars and trucks. But forced to feed so many brands, G.M. often resorted to a practice called
``launch and leave'' -- spending billions upfront to bring vehicles to market, but then failing to
keep supporting them with sustained advertising.

With its market share shrinking, G.M. could not give its multiple brands and car models the
individual attention that helped Honda attract customers to the Accord and Toyota to its Camry.

It also lost interest in vehicles that needed time to find their audience, as happened when the
company introduced the EV1 electric vehicle and then dropped it in 1999 after only three years.

Now G.M.'s brand lineup is being halved, with the company jettisoning divisions like Pontiac.

``Nobody gave any respect to this thing called image because it wasn't in the business plan,''
Mr.~Wangers said. ``It was all about, 'When is this going to earn a profit?' ''

Over the years, G.M. executives became practiced at the art of explaining their problems,
attributing blame to everyone but themselves.

That list included the United Automobile Workers, for demanding health care coverage and pensions
(even though G.M. agreed to provide them); government regulators, for imposing rules that G.M. said
hampered its competitiveness; the Japanese government, for unfairly helping its own carmakers break
into the United States market; and the news media, for failing to appreciate G.M. vehicles and the
strides the company was making to improve them.

Asked in 1995 why he had not moved faster to reorganize the company, the late G.M. chief executive
Roger Smith replied, ``Wouldn't it have been wonderful if we could have flipped a switch?''

Even last week, G.M.'s newly retired vice chairman, Robert A.~Lutz, said the automaker had
experienced a ``world of hurt, much of it not of our own doing.''

Sloganeering was not backed up by execution. Executives wore lapel pins, for example, in 2002, with
the number ``29'' -- referring to the market share the company vowed to regain (most companies focus
on profits). Through April of this year, its share was 19 percent, a steep drop from its peak of 54
percent in 1954.

Consumers started blaming G.M. for sub-par vehicles. They may have given them second and third
chances, but many eventually started switching to other brands, which will make it that much harder
for G.M. to win them back.

Mr.~Wagoner was able to hold on to his job for longer than people expected, as G.M.'s stock fell
steadily from about \$70 when he took charge at the start of the decade. It closed at 75 cents a
share on Friday.

Mr.~Wagoner was pushed out by the Obama administration, which is now making the call to push the
company into bankruptcy court.

A judge will then start the process of building a new, though much diminished, G.M. into a company
that might have a shot at a second century. But the automaker that so dominated center stage in the
American car market for so long will have to earn that place back.


\section{G.M. to Seek Bankruptcy and a New Start}

\textsl{This article was reported by David E. Sanger, Jeff Zeleny and Bill Vlasic, and written by
  Mr. Sanger.}

WASHINGTON -- President Obama will push General Motors into bank-ruptcy protection on Monday, making
a risky bet that by temporarily nationalizing the onetime icon of American capitalism, he can save
at least a diminished automaker that is competitive.

The bankruptcy, to be filed in New York, is a moment of reckoning for an industry that was once at
the heart of the American economy. It culminates a remarkable four months of confrontation between
Washington and Detroit that is expected to result in a drastic downsizing of the company.

It also places the government in uncharted territory as a business owner, as it takes a 60 percent
ownership stake in the company during its restructuring.

Reflecting the government's extraordinary intervention in industry, aides say, Mr.~Obama plans to
tell the nation on Monday that he believes G.M. can be brought back from the brink of insolvency,
even if the company looks almost nothing like the titan of old.

Meanwhile, a federal judge late Sunday night cleared a path for Chrysler to get out of bankruptcy by
approving a sale of most of that carmaker's assets to a new entity to be run by Fiat of Italy.

Administration officials briefed reporters on the G.M. plans Sunday night, as President Obama began
to inform members of Congress. But the White House insisted that the aides who talked to reporters
could not be named.

In his remarks on Monday, Mr.~Obama will spell out a strategy in which a shrunken G.M. can make
money even if new car sales remain at a sluggish 10 million a year in the United States and even if
G.M., once the giant of the industry, drops below its current 20 percent market share in this
country.

But to get there, American taxpayers will invest an additional \$30 billion in the company, atop
\$20 billion already spent just to keep it solvent as the company bled cash as quickly as Washington
could inject it. Whether that investment will ever be recovered is still an open question.

The company will also have to shed 21,000 union workers and close 12 to 20 factories, steps that
most analysts thought could never be pushed through by a Democratic president allied with organized
labor.

Forty percent of the company's 6,000 dealers will close, the workers' union will be forced to
finance half of its \$20 billion health care fund with stock of uncertain value in the restructured
G.M., and bondholders, including many retirees, will be forced to take stock worth 10 cents for
every dollar they lent the company.

The company's last steps toward bankruptcy took place over the weekend as a majority of G.M.
bondholders agreed not to challenge the filing in court and to exchange their debt for stock.

Lawrence H.~Summers, who as head of the National Economic Council serves as one of the co-heads of
the auto task force, argued in an interview on Sunday that the bailout of the auto industry was
fundamentally different from the Mexican bailout in 1994, the Asian economic crisis in the late
1990s, and the continuing banking crisis.

General Motors and Chrysler, he said, were ``clear cases of insolvency,'' in which mere loans would
not accomplish the goal of getting the automakers past a temporary crisis. ``There was no argument
that they were solvent, no argument they could meet their obligations.''

He said that left the Obama administration to decide whether to allow ``a laissez-faire,
uncontrolled bankruptcy, which would have had an enormous cost,'' or a ``controlled process,'' in
which the goal was to make sure that the auto companies not only restructured, but were not
overburdened with debt. So, in return for what amounted to debtor-in-possession financing, Mr.~Obama
chose to accept equity in the new company -- while insisting that he had no intention of exercising
day-to-day control over the company.

``It's a fine line,'' Mr.~Summers said, ``but we think it is manageable.''

To assist in the restructuring, the automaker is expected to hire the consulting firm Alix Partners,
which has worked on several major bankruptcies, including those for Enron and Kmart. One of the
firm's partners, Al Koch, is expected to manage the liquidation of corporate assets that G.M. will
shed during its Chapter 11 restructuring, people with knowledge of the strategy said.

Mr.~Obama is taking several risks under the plan. None may be bigger than the decision that the
United States government will take its 60 percent share of the stock in a new G.M., leaving
taxpayers vulnerable if the overhaul is not successful. (Canada, for its part, is taking a 12
percent stake.)

``We don't think that after this next \$30 billion, they will need more money,'' one administration
official said. ``But the fact is there are things you don't know -- like when the car market will
come back, and how much Toyota and Honda and Volkswagen will benefit from the chaos.''

On Monday, Mr.~Obama is expected to argue that any alternative to his plan would be worse, and that
a liquidation of G.M. -- the only other real option -- would send the unemployment rate soaring over
10 percent and would radiate damage throughout the economy.

But aware of the hardships the plan will impose on regions across the country that depend on auto
production, the White House is dispatching a dozen Cabinet members and other officials across four
states this week to reassure residents.

Aides say Mr.~Obama will portray himself on Monday as a reluctant shareholder, eager to sell the
company back to private investors, perhaps within 6 to 18 months.

Officials say the president will insist that once the government sets up new management and a board
of directors, it will remove itself from G.M.'s day-to-day operations. But even his aides anticipate
intense pressure as the company's managers are called to testify in Congress and face questions like
why they decided to build new cars in Mexico and South Korea, rather than in Michigan or the South.

``Congress and many Americans are going to say, if we own it, why can't we make these decisions?''
one of Mr.~Obama's top economic aides said, ``and it's going to be a challenge to answer that.''

To ease the way, the White House on Sunday briefed reporters on a new set of principles for how the
government should behave as a majority shareholder. It argued that the government's role should be
limited primarily to the beginning of the process, but that it should then recede, becoming a
passive investor, one seeking to sell its stake quickly.

At the same time, Mr.~Obama has laid out goals for all the Detroit automakers that will presumably
affect their major strategic decisions. He has urged them, for example, to build smaller cars with
significantly better fuel efficiency.

Six months ago, even the suggestion of such deep intervention into G.M.'s operations would have
raised huge objections. But by the time the denouement came, the company seemed almost relieved.
Robert Lutz, G.M.'s vice chairman, said that ``for the first time in our history, the American auto
industry has the ear of the administration. Their number one goal is to make us successful.''

Nonetheless, Michael Useem, a professor of management at the Wharton School at the University of
Pennsylvania, said the decision would ``mean a new chapter in the history books on American
capitalism.'' He added, ``How we think about American free enterprise is really hanging in the
balance.''

For Mr.~Obama, whose ascent to the White House depended on carrying states across the industrial
Midwest, the political risk is significant.

The G.M. bankruptcy will ripple across several states where hundreds of parts suppliers and car
dealerships face imminent closings.

Indeed, the four states where Cabinet secretaries are focusing their efforts this week -- Indiana,
Michigan, Ohio, Wisconsin -- all were carried by Mr.~Obama last November. It was the first time
Indiana has supported a Democratic presidential candidate in 44 years.

These Main Street political challenges will almost certainly be an issue for Democrats on the ballot
in next year's midterm election campaign and in the president's own re-election effort in 2012. If
those jobs shift to nonunion plants in the South, where German and Japanese carmakers have built
their facilities, or overseas, Mr.~Obama could face criticism inside his own party.

``It is unacceptable to ask U.S.~workers to subsidize the exportation of their own jobs,'' said
Representative Dennis Kucinich, Democrat of Ohio, whose district includes Cleveland. ``The
taxpayers' investment should be used to protect American plants so that American workers can build
the next generation of automobiles.''

In his presidential campaign speeches last year, often delivered in the shadow of closed
manufacturing plants, Mr.~Obama bluntly conceded that most of the jobs would not come back. Instead,
his administration is pointing to investments that the economic recovery act will make in
communities.

Rob McNabney, chairman of the Madison County Democratic Party in Anderson, Ind., a onetime booming
automotive center, said the problems for Mr.~Obama were severe. ``He's going to be judged by what he
does,'' Mr.~McNabney said.

\end{document}
